\subsection*{等待的是竞争事务的结局}

这里继续采用前面说明的 PostgreSQL 18 Read Committed。普通查询的语句快照，与数据库为了检查唯一性而等待竞争事务，是两项相互配合的规则。

假定 A 已插入 S，但尚未提交。B 的普通 SELECT 查不到 A 的未提交记录，B 随后也尝试插入 S。数据库检查唯一约束时，必须考虑可能提交的竞争记录。如果忽略 A，允许 B 也提交，就可能产生两条相同唯一键。因此这项检查会等待 A 的事务结局；它承担的检查范围与普通 SELECT 的数据可见范围有区别。

PostgreSQL 18 手册对唯一性检查的说明给出了这条规则：发现相同键由未结束的其他事务插入时，要等待该事务完成，再重新判断唯一性。在 Read Committed 中，DO NOTHING 可能因另一个事务的结果而跳过插入，即使那条竞争记录对当前插入语句的快照仍不可见\cite{pg}。两条规则共同决定后续的等待与读取。

\begin{figure}[htbp]
\centering
\begin{tikzpicture}[x=1cm,y=1cm]
\node[booknote] at (2,0.5) {A：先插入 S};
\node[booknote] at (8,0.5) {B：尝试同一 S};
\node[booknode,text width=3.8cm] (a1) at (2,-.3) {S → P1\\尚未提交};
\node[bookresult,text width=4.5cm] (b1) at (8,-.3) {唯一性检查发现竞争\\等待 A 的事务结局};
\node[booknode,text width=3.8cm] (a2) at (2,-1.8) {扣款、订单、流水\\保存原结果并提交};
\node[booknode,text width=4.5cm] (b2) at (8,-1.8) {A 提交后重新检查\\DO NOTHING，返回零行};
\node[booknode,text width=4.5cm] (b3) at (8,-3.3) {下一条 SELECT 读取 P1\\同参，返回原成功};
\draw[bookarrow] (a1)--(a2);
\draw[bookarrow] (b1)--(b2);
\draw[bookarrow] (a2.east)--node[above,font=\footnotesize]{提交后唤醒}(b2.west);
\draw[bookarrow] (b2)--(b3);
\node[booknote,text width=3.3cm] at (12.5,-1.8) {B 全程\\未进入扣款};
\node[booknote,text width=13cm] at (7,-4.45) {最终：S → P1，只有 O1、D1，账户余额 7000};
\end{tikzpicture}
\caption{竞争者提交以后，等待者跳过插入，再以新语句读取原结果。}
\label{fig:ch04-winner-commit}
\end{figure}


从初始余额 10000、购买表为空开始，A 插入 P1 对应的 S，取得继续处理的资格。B 插入同一 S 时停在等待，尚未进入扣款。A 扣款并写 O1、D1，最后提交。B 被唤醒，重新检查时发现 S 的唯一记录已经成立，DO NOTHING 跳过它的候选 P2，RETURNING 返回零行。

B 接着发出另一条 SELECT，读取 P1 的规范化版本、输入、终态和原结果。两份输入相同，它返回 P1、O1、D1 及 7000。最终只有一条成功扣款，余额保持 7000。B 的候选 P2 没有被插入，也没有对应订单。服务从查到的 P1 构造返回内容，客户端于是可以用 P1 继续查询实际成立的购买。

数据库监控可以把这类等待显示为 Lock 类型的 \code{transactionid} 等待，用来定位尚未结束的竞争事务。应用仍以插入返回行数选择分支，随后读取终态；诊断中的等待名称不代替这两项业务判断。

\subsection*{原竞争者回滚以后谁继续}

回到相同初始数据，A 插入 S 后还没提交，B 继续在同键插入处等待。A 写订单时遇到错误，整个事务回滚。A 曾修改的余额恢复，P1、O1、D1 和它对 S 的占用都没有成为已提交事实。B 被唤醒，重新检查发现阻挡它的键已不再成立，便插入自己的候选 P2；RETURNING 这次返回 P2。

\begin{figure}[htbp]
\centering
\begin{tikzpicture}[x=1cm,y=1cm]
\node[booknote] at (2,0.5) {A：先插入 S};
\node[booknote] at (9,0.5) {B：尝试同一 S};
\node[booknode,text width=4.2cm] (a1) at (2.5,-.3) {S → P1\\事务内部可能已经扣款};
\node[bookresult,text width=4.8cm] (b1) at (10,-.3) {等待 A 结束\\尚未取得身份};
\node[booknode,text width=4.2cm] (a2) at (2.5,-1.9) {A 出错，整笔回滚\\S 占用与扣款一起撤销};
\node[booknode,text width=4.8cm] (b2) at (10,-1.9) {重新检查，插入 P2 成功\\RETURNING 返回 P2};
\node[booknode,text width=4.8cm] (b3) at (10,-3.5) {继续报价与条件扣款\\完成终态后提交};
\draw[bookarrow] (a1)--(a2);
\draw[bookarrow] (b1)--(b2);
\draw[bookarrow] (a2.east)--node[above,font=\footnotesize]{回滚后唤醒}(b2.west);
\draw[bookarrow] (b2)--(b3);
\node[booknote,text width=5cm] at (3,-3.5) {P1 没有已提交终态\\B 在自己的事务内能看到 P2 的 pending};
\node[booknote,text width=13cm] at (7,-4.8) {若 B 成功，最终 S → P2、O2、D2，余额 7000；用户操作仍为同一 S};
\end{tikzpicture}
\caption{原事务回滚以后，等待者才有机会接着完成原操作。}
\label{fig:ch04-winner-rollback}
\end{figure}


B 取得的是继续尝试购买的资格。它仍须执行报价与余额条件，写入完整终态并提交。若这些步骤成功，余额由 10000 变为 7000，最终保留 P2、O2、D2；P1 没有已提交结果。P2 的数字后缀只帮助区分这两次候选，用户操作身份始终是同一 S，这次恢复完成的仍是 Alice 原先确认的购买。

若 A 回滚之前只完成身份插入，B 的后续路径相同；若 A 已在事务内部扣过款，回滚会一起撤销扣款。这个结论依赖修改都属于同一事务，并且没有在其中调用一个已经独立完成扣款的外部系统。应用日志即使记下候选编号 P1，也只说明该尝试走到了编号分配步骤；P1 是否成交，仍要以已提交的购买、订单和扣款记录核对。

B 插入 P2 后，在自己的事务里可以读到 P2 的 \code{pending}；其他连接的普通查询仍读不到它。若 B 填齐成交快照、订单、流水与终态并提交，P2 就成为外部可见的成功购买；若 B 在提交前也回滚，独立连接最终查不到 S，余额恢复为 10000。事务内部看见占位，与外部已经形成终态，需要按这两个观察范围区分。

这里可以把提交和回滚两条分支按同一组观察对象比较。A 未结束时，B 只知道有竞争，暂时没有原业务结论。A 提交成功以后，B 才有 P1 的稳定结果可以返回。A 回滚以后，B 也没有一个可以重放的 P1 拒绝结果，只有重新尝试的机会。事务回滚说明本次尝试的数据库修改撤销，业务上的余额不足则是可以提交的明确结论；恢复流程据此选择继续尝试或返回终态。

\begin{figure}[htbp]
\centering
\begin{tikzpicture}[x=1cm,y=1cm]
\node[booknote] at (2,0.45) {竞争事务的结局};
\node[booknote] at (7,0.45) {持久事实};
\node[booknote] at (12,0.45) {等待者下一步};
\node[booknode,text width=3.7cm] (ok) at (2,-.45) {提交 succeeded};
\node[booknode,text width=3.7cm] (oy) at (7,-.45) {S、P1、O1、D1\\原成功结果};
\node[booknode,text width=3.7cm] (or) at (12,-.45) {同参返回原成功};
\node[booknode,text width=3.7cm] (no) at (2,-2) {提交余额不足};
\node[booknode,text width=3.7cm] (ny) at (7,-2) {S、P1、原拒绝\\订单与扣款为空};
\node[booknode,text width=3.7cm] (nr) at (12,-2) {同参返回原拒绝};
\node[booknode,text width=3.7cm] (ab) at (2,-3.55) {事务回滚};
\node[bookresult,text width=3.7cm] (ay) at (7,-3.55) {本次尝试无已提交终态\\占用与修改撤销};
\node[booknode,text width=3.7cm] (ar) at (12,-3.55) {重新竞争 S\\按实际条件继续};
\draw[bookarrow] (ok)--(oy);\draw[bookarrow] (oy)--(or);
\draw[bookarrow] (no)--(ny);\draw[bookarrow] (ny)--(nr);
\draw[bookarrow] (ab)--(ay);\draw[bookarrow] (ay)--(ar);
\end{tikzpicture}
\caption{业务拒绝作为终态提交，事务回滚则把本次尝试的修改撤销。}
\label{fig:ch04-outcome-evidence}
\end{figure}


例如 A 的扣款语句返回零行以后，服务正常保存 \code{insufficient\_funds} 并提交。B 等待后读到这条终态，就跟随原结论。A 的扣款语句尚未完成时进程退出，数据库回滚，B 取得 S 后重新检查余额，则可能得到另一个结果。后一历史中，此前没有终态提交，所以 B 依据自己实际处理时的余额决定购买结果。

\subsection*{一条语句的空结果可能跨过一次提交}

可以想到用一条语句同时完成插入和原记录查找，减少服务与数据库的一次往返。常见写法使用 WITH，把插入放进一个有名称的中间结果，再把新插入行和已有行合并。WITH 引入的这种命名子查询称为公共表表达式；它让同一条 SQL 的其他部分引用该计算结果，本例把它命名为 \code{inserted}。

\begin{BookCode}
WITH inserted AS (
  INSERT INTO purchases (...)
  VALUES (...)
  ON CONFLICT (account_id, action, session_id, idempotency_key)
  DO NOTHING
  RETURNING purchase_id
)
SELECT purchase_id FROM inserted
UNION ALL
SELECT purchase_id FROM purchases
WHERE account_id = 1
  AND action = 'purchase.create.v1'
  AND session_id = 'E1'
  AND idempotency_key = 'K1'
LIMIT 1;
\end{BookCode}

省略号仅代表已经列出的完整插入字段，这个片段用于分析快照。UNION ALL 把两个查询结果合在一起，LIMIT 1 只取其中一行。设这条语句开始时 A 仍未提交，B 的语句快照尚看不到 P1。插入部分检测到相同 S，于是等 A。A 随后提交，B 的唯一性检查据此放弃插入，\code{inserted} 为空。

剩下的 SELECT 仍属于同一条语句，继续使用语句开始时的可见范围。A 在等待期间提交的 P1，对这个快照仍不可见，于是第二个分支也为空，整条语句返回零行。数据库此刻已经有 P1，B 也确实因为 P1 而没有插入，但这两点没有使同语句的普通 SELECT 换成更新快照。

\begin{figure}[htbp]
\centering
\begin{tikzpicture}[x=1cm,y=1cm]
\node[booknode,text width=4.7cm] (snap) at (3,0) {B 的语句开始\\快照中没有未提交的 P1};
\node[booknode,text width=4.7cm] (ins) at (11,0) {WITH 内插入同一 S\\唯一性检查等待 A};
\draw[bookarrow] (snap)--(ins);
\node[bookresult,text width=4.7cm] (commit) at (3,-1.7) {A 在等待期间提交 P1\\唯一检查据此判断冲突};
\node[booknode,text width=4.7cm] (skip) at (11,-1.7) {插入被跳过\\inserted 为空};
\draw[bookarrow] (ins.south)--(skip.north);
\draw[bookalert] (commit)--(skip);
\node[booknode,text width=4.7cm] (select) at (3,-3.4) {同语句 SELECT 沿用旧快照\\已有行查询也为空};
\node[booknode,text width=4.7cm] (next) at (11,-3.4) {B 发出下一条 SELECT\\新快照读到 P1};
\draw[bookarrow] (skip.south)--++(0,-.4)-|(select.north);
\draw[bookarrow] (select)--node[above,font=\footnotesize]{新语句}(next);
\node[booknote,text width=13cm] at (7,-4.8) {唯一检查可以等待事务结局；本条普通 SELECT 的快照在等待后保持不变};
\end{tikzpicture}
\caption{同一语句可能因刚提交的行放弃插入，却读不到那条行。}
\label{fig:ch04-statement-snapshot}
\end{figure}


应用若把这个空结果解释成没有原操作，再绕开唯一约束扣款，就会重新制造失败路径。此时需要更新读取范围，沿 S 找到刚才阻止插入的已提交事实。

本章采用两条语句：插入结束后，冲突分支再发 SELECT。Read Committed 下，新语句取得新的读取视图，能看到前面等待后已经提交的 P1。两条语句可以位于同一事务中。这个行为依赖隔离级别按语句更新读取视图；把相同代码移到另一隔离级别以后，需要重新检查。

在记录保持不删、读取同一主库、按精确 S 查找且没有额外过滤的本章路径中，竞争提交后的下一条 SELECT 应当找到原记录。实现仍要明确处理找不到的分支：记录若被其他维护流程删除，或读取路径与预期不同，就结束当前尝试后按原身份恢复，或者在有界次数内重新走唯一插入流程。循环指重复执行同一套受约束步骤，每一轮仍以实际插入成功作为进入购买的条件。达到次数或等待上限时，返回可恢复的暂时错误，保留 S，并查明记录保留或读取路径为何与预期不同。

\subsection*{等到原记录以后还要核对输入}

让 A 使用 S 和 \code{paper-001}、数量 1、Q7，B 错把同一 S 的报价换为 Q8。沿用第三章的报价关系，Q7 为 3000 分，Q8 为同一商品的新报价 3500 分。两个请求的身份相同，输入不同。B 插入时依然会等待 A，唯一约束比较的是 S，所以输入差异不会给 B 另开一条同身份记录。

A 提交以后，B 的插入返回零行。B 按保存版本规范化自己这次输入，与 P1 的 \code{canonical\_input} 比较，得到不同，返回 409 \code{conflict}。P1 仍然表示按 Q7 完成的购买，原输入保持不变，余额仍为 7000；B 携带 Q8 的这次尝试得到输入冲突。

\begin{figure}[htbp]
\centering
\begin{tikzpicture}[x=1cm,y=1cm]
\node[booknode,text width=4.3cm] (a) at (3,0) {A：S、paper-001、Q7\\插入并完成 P1};
\node[booknode,text width=4.3cm] (b) at (11,0) {B：S、paper-001、Q8\\在同 S 插入处等待};
\node[booknode,text width=4.3cm] (old) at (3,-1.8) {A 提交后的绑定\\S → 原输入与 P1};
\node[bookresult,text width=4.3cm] (check) at (11,-1.8) {B 插入零行，再读取\\按保存版本比较输入};
\draw[bookarrow] (a)--(old);\draw[bookarrow] (b)--(check);
\draw[bookarrow] (old)--(check);
\node[bookresult,text width=10cm] (conflict) at (7,-3.5) {报价不同：409 conflict\\P1 仍按 Q7 成交；余额保持 7000};
\draw[bookalert] (check.south)--(conflict.north);
\end{tikzpicture}
\caption{同身份异参先经过同一唯一竞争，再按已经成立的绑定报告冲突。}
\label{fig:ch04-input-conflict}
\end{figure}


若 A 回滚，B 可能插入成功；Q8 仍有效且账户足够支付时，B 按 3500 首次成交，余额为 6500，并将 S 绑定到 Q8。A 以后拿 Q7 恢复，就会收到异参冲突。这个设计采用首次成功提交的绑定，会话许可没有预先固定商品与报价。

协议要求同一操作的所有尝试从同一份持久输入构造；把同一 S 配给 Q7 和 Q8，违反了调用方的输入稳定规则。即使遇到这种输入混用，服务仍须保证：S 至多按一份完整输入提交，金额与该输入相符；已提交绑定不可覆盖；另一份输入得到明确冲突。服务仅从两份相互矛盾的输入中无法判断哪份符合用户本意。需要更早固定输入的系统，可以在许可签发时保存摘要和完整解释依据，并在所有入口检查，同时承担预登记阶段的状态与恢复规则。

相反，\(S_1=(1,\texttt{purchase.create.v1},E1,K1)\) 与 \(S_2=(1,\texttt{purchase.create.v1},E1,K2)\) 是两个完整身份。它们都携带 \code{paper-001}、数量 1、Q7 时，各自可以插入成功，随后在账户行的更新上排队。初始 10000 足够两笔各扣 3000，最终余额 4000，两笔购买都合法。若给规范化输入单独加唯一约束，Alice 明确确认的两份相同商品就会被合并，因此本章按身份约束受理次数，按输入核对原绑定。

\subsection*{同账户的新操作怎样检查余额}

唯一身份解决重复受理，同账户的不同购买仍共享余额。以两条不同 S、初始余额 5000、每笔扣 3000 为例，两笔都可以占住各自身份。若 A 先更新账户，把内部余额改为 2000 而尚未提交，B 的条件 UPDATE 会等待 A 对该行的修改结束。

在 PostgreSQL 18 Read Committed 下，A 提交后，B 尝试在已更新版本上应用 UPDATE，并重新判断 WHERE 条件。新版本余额为 2000，\code{balance\_minor >= 3000} 不成立，B 更新零行。B 把自身购买保存为 \code{insufficient\_funds}，提交原拒绝结果；最终成功购买一笔，拒绝一笔，余额为 2000。若 A 回滚，B 可基于恢复后的 5000 成功扣款。

\begin{figure}[htbp]
\centering
\begin{tikzpicture}[x=1cm,y=1cm]
\node[booknote,text width=13cm] at (7,.45) {初始余额 5000；S1 与 S2 不同，各申请扣 3000};
\node[booknode,text width=4.8cm] (a) at (3,-.6) {A 取得 S1\\账户更新为 2000，尚未提交};
\node[booknode,text width=4.8cm] (b) at (11,-.6) {B 取得 S2\\UPDATE 等待账户行};
\node[booknode,text width=4.8cm] (ac) at (3,-2.3) {A 提交成功\\余额 2000};
\node[bookresult,text width=4.8cm] (bc) at (11,-2.3) {B 按新版本重查余额条件\\$2000<3000$，返回零行};
\draw[bookarrow] (a)--(ac);\draw[bookarrow] (b)--(bc);
\draw[bookarrow] (ac)--node[above,font=\footnotesize]{提交后继续}(bc);
\node[booknode,text width=10cm] at (7,-4) {S1：成功，订单与扣款各一条\\S2：余额不足，保存原拒绝；最终余额 2000};
\end{tikzpicture}
\caption{不同身份分别受理，再通过账户条件更新决定哪笔资金变化成立。}
\label{fig:ch04-balance-contention}
\end{figure}


这条重新检查条件的行为由 PostgreSQL 18 的隔离级别说明给出\cite{pg}。它适用于这里的单行条件更新。若业务条件改为跨多个账户或多条订单累计判断，分别读取多行再写入，需要重新设计锁定、约束或隔离级别。

把余额读到应用里，再用旧值计算固定余额，执行 \code{UPDATE ... SET balance\_minor=2000}，会遗漏上述基于新版本余额重算的业务关系。原 SQL 的减法与余额条件在数据库处理该行时一起判断，使两个不同 S 竞争账户时的失败结论明确。身份约束与条件更新分别观察不同对象，两项规则都要保留。

还可用 10000 分演算两项不同 S 同时扣 3000 的情况。A 更新后内部余额 7000，B 等待；A 提交以后，B 在新版本 7000 上重新检查，条件成立，减法得到 4000，B 返回 4000 并提交。两笔流水的当次余额依次为 7000 与 4000。此时 4000 是两项合法购买的结果，与第一节同一 S 重复受理得到的 4000 数值相同，业务含义由身份关系区分。

账户查询只有余额时，看不出哪条历史发生。核查需要沿购买身份取出对应订单和流水，并确认 S1、S2 是否来自两次明确确认。单一最终数值既可能来自正确执行，也可能来自重复操作，因此演算同时保留余额、购买与输入。

条件 UPDATE 返回零行时，如果 WHERE 从快照中就排除了账户行，事务可能没有取得该账户的行锁。因此，这个分支不自动保证账户保持原值直到拒绝提交。第 6.2 节会以购买扣款与并发充值两种动作集合，分别检查余额不足可以放在哪个生效位置；这里先保留单行条件判断与完整终态共同提交的规则。

\subsection*{事务错误以后从哪里恢复}

如果使用普通 INSERT，重复 S 会引发唯一冲突错误。在 PostgreSQL 中，一条语句使事务进入失败状态后，继续发普通 SELECT 通常得到当前事务已中止的错误。服务需要先回滚整个事务，再开始读取；也可以预设保存点，把局部失败回滚到该位置。保存点是事务内部可返回的中间位置，用于撤销其后部分工作。本章用指定冲突目标的 DO NOTHING，让预期身份冲突走正常返回分支，减少这条局部错误路径。

DO NOTHING 只处理指定的唯一冲突，其他 SQL 错误仍可能使事务失败。锁等待超时也需要按数据库错误和连接状态回滚处理。锁等待超时表示某次等待超过允许时长，等待者仍未取得身份，竞争者也可能继续运行。返回暂时不可用时，客户端继续保存 S；换成 K2 再发，会把原来仍可能提交的 K1 旁边增加一项新购买。

死锁是两个或多个事务各自占有对方还需要的资源，形成相互等待的环。例如一个事务先锁 S1 再锁账户，另一个先锁账户再尝试处理 S1，就可能形成环。数据库会检测并终止其中一个事务，让另一方继续。统一资源取得次序有助于减少死锁；被终止者重新开始整笔事务时，仍沿用原账户、动作、E1、K 和已确认输入。原事务已明确回滚，它的新尝试还要重新竞争 S，其他同身份请求可能已经接着完成。

Repeatable Read 是 PostgreSQL 的另一隔离级别，即可重复读。它让事务内后续普通读取保持在同一事务快照范围，避免同一事务反复读取时看到其他事务新提交的变化。于是本章依赖下一条 SELECT 获得新可见范围的代码，需要按该级别重新安排；并发修改还可能产生需要回滚重试的序列化错误。隔离级别改变读取与错误规则，客户端仍需要稳定操作身份。

Serializable 通常称为可串行化，要求成功提交的并发事务在数据库效果上可以解释为某个逐笔执行顺序。PostgreSQL 用依赖检测等机制发现与这个要求不相容的并发历史，并可能取消一笔事务，要求应用整体重试\cite{pg}。序列化失败明确发生时，应用从事务开头重新执行读取、判断和写入；所有尝试继续使用原 S。若根本没有保存 S，不同事务即使被排成先 P1 后 P2 的顺序，仍然可能各扣 3000，这个顺序自身满足数据库可串行化，却不符合 Alice 只授权一次的要求。

还要区分明确回滚与提交确认丢失。服务发送 COMMIT 后连接中断，应用可能无法知道数据库是否已经提交。它应结束失效连接，用原 S 重新进入受约束流程：原事务提交，就读取原结果；原事务回滚，就有机会重新受理；原事务仍未结束，竞争可能继续等待。这里还依赖原数据库事实没有丢失、身份记录仍保留和恢复读取路径正确。连接错误留下的是结果待确认，恢复流程要继续核查同一身份。

\begin{figure}[htbp]
\centering
\begin{tikzpicture}[x=1cm,y=1cm]
\node[bookresult,text width=10cm] (lost) at (7,0) {COMMIT 发出后连接中断\\本次尝试暂时无法确认结果};
\node[booknode,text width=3.7cm] (committed) at (2,-1.8) {原事务已提交\\S 有原终态};
\node[booknode,text width=3.7cm] (aborted) at (7,-1.8) {原事务已回滚\\S 尚无原终态};
\node[booknode,text width=3.7cm] (running) at (12,-1.8) {原事务仍进行\\S 存在竞争};
\draw[bookarrow] (lost.south west)--(committed.north);
\draw[bookarrow] (lost.south)--(aborted.north);
\draw[bookarrow] (lost.south east)--(running.north);
\node[booknode,text width=3.7cm] (r1) at (2,-3.4) {按原 S 读取\\核对并返回原结果};
\node[booknode,text width=3.7cm] (r2) at (7,-3.4) {按原 S 再次竞争\\可能接着完成};
\node[booknode,text width=3.7cm] (r3) at (12,-3.4) {按原 S 尝试插入\\等待或暂时结束尝试};
\draw[bookarrow] (committed)--(r1);\draw[bookarrow] (aborted)--(r2);\draw[bookarrow] (running)--(r3);
\node[booknote,text width=13cm] at (7,-4.6) {三个分支都沿用原账户、动作、E1、K 与输入；连接错误没有创建新购买授权};
\end{tikzpicture}
\caption{提交确认丢失以后，以原身份恢复才能分辨尚未确定的事务结局。}
\label{fig:ch04-recovery-outcomes}
\end{figure}


\subsection*{把一次等待安排到完整服务中}

服务线程等待数据库时，客户端自己的等待期限可能先到。假设 A 仍在处理，B 已经等待相同 S 两秒，客户端两秒后显示结果待确认。这只说明客户端停止等待 B 的响应，A 与 B 的数据库操作仍可能继续。服务可以设定较短的数据库等待上限和整体请求上限，减少长时间占用连接；达到上限时的暂时错误仍要保留原操作的恢复入口。

如果 B 的数据库语句被明确取消并整笔回滚，B 这次没有提交购买，A 仍可能提交。客户端以后继续原 S，依然可能读到 A 的成功结果。若 B 在客户端断开后继续等待，最终读到 P1，即便它无法把响应交给原连接，P1 的结果也仍然保存在数据库。恢复者再来查找，可以用新的通信尝试读取这份结果。

服务可以把同一进程中正在等待同 S 的请求合并，减少数据库连接压力；所有参与者等待同一份结果或一起转入恢复。这里的合并是一种资源优化，需要处理进程退出和订阅者断开，持久唯一约束仍留在数据库。若优化层丢失状态，后来的请求照常走受约束插入，再根据数据库的实际返回选择分支。

一次数据库事务也应尽量短。读取本地报价快照、执行条件更新和填写结果都应有明确边界；等待用户确认或调用不可控的远程服务，会延长持有 S 与账户锁的时间。远程动作还可能产生本地回滚覆盖不到的效果。本章只使用已经可以在当前事务内验证的 Q7 依据，账户扣款由本地条件 UPDATE 完成。
\FloatBarrier
